\section{Inputs, normalisation and computation}\label{sec:inputs}

\subsection{Torus knots}
For the torus knot $T_{m,n}$ we compute $\brCS_\lambda$ from the action of the generators of the elliptic Hall algebra, following Gorsky and Negu\c{t} \cite{Gorsky13b}.
The relation between the elliptic Hall algebra and the DAHA is due to Schiffmann and Vasserot \cite{Schiffmann09}, and its use for cables to Morton and Samuelson \cite{Morton14}.
% [I1]
To fix conventions we reproduced, coefficient by coefficient, the $27$ values for torus knots in the ancillary files of KN
($T_{2,3}$, $T_{2,5}$, $T_{2,7}$, $T_{2,9}$ with $|\lambda|\le3$ and $T_{3,4}$ with $\yng(1),\yng(2),\yng(1,1)$)
and the $13$ tables of BPS invariants in Appendix~C of KN, under the monomial normalisation of \S\ref{sec:normalization}.
This is a reproduction of known results and is used only as a check of conventions.

\subsection{Algebraic iterated torus knots}
We write $C_{m,n}K$ for the $(m,n)$-cable of a companion knot $K$; $m$ is the winding number and $n$ is measured in the Seifert framing of the companion.
In the DAHA definition of Cherednik and Danilenko \cite{Cherednik14,Cherednik15} the knot $C_{m_2,n_2}T_{m_1,n_1}$ has Newton pairs $r=(m_1,m_2)$, $s=(n_1,n_2-m_1n_1m_2)$,
and it is algebraic (the link of a plane curve singularity) when all $r_i,s_i>0$, i.e.\ $n_2>m_1n_1m_2$.
% [I2]
We decompose this definition linearly at each step of the iteration and compute it as a composition of operators of the elliptic Hall algebra.
For the outermost step $C_{m,n}$ we expand
$\Phi_\lambda=P_\lambda[p_k\mapsto P_{k(m,-s')}]\cdot1=\sum_\mu c_{\lambda\mu}(s')P_\mu$,
where $s'$ is the last component of $s$, $p_k$ are power sums and $P_{k(m,-s')}$ are the generators of slope $(m,-s')$,
and take the linear combination with the values $\rCS_\mu(K)$ of the companion.
It is known that the naive $t$-deformation of the Rosso--Jones formula is incorrect \cite{Danilenko14,Morozov18}; the form used here follows directly from the definition.
\begin{compresult}\label{cr:daha}
The values so computed agree, in all coefficients and up to a monomial, with the ten explicit examples of Cherednik and Danilenko:
$K1$ with $\yng(1)$ and $\yng(2)$, and $K2,K3,K4,K5a,K5b,K5c,K7,K8$ with $\yng(1)$ (the non-algebraic cables are named in Appendix~\ref{app:nonalg}).
At $t=q$ they agree with the coloured HOMFLY-PT polynomials computed from braid presentations.
\end{compresult}
The cables studied in this paper are listed in Table~\ref{tab:cables}.
\begin{table}[h]
\centering
\caption{Algebraic iterated torus knots. $b$ is the number of strands of the braid and $\delta$ the genus-type quantity of Computational Result~\ref{cr:M}.}\label{tab:cables}
\begin{tabular}{llcc}
\toprule
Name & Knot & $b$ & $\delta$ \\
\midrule
$K1$ & $(2,13)$-cable of $T_{2,3}$ & 4 & 8 \\
$K3$ & $(2,15)$-cable of $T_{2,3}$ & 4 & 9 \\
$K6a$ & $(2,21)$-cable of $T_{2,5}$ & 4 & 14 \\
$K9$ & $(2,31)$-cable of $T_{2,7}$ & 4 & 21 \\
$K8$ & $(2,25)$-cable of $T_{3,4}$ & 6 & 18 \\
$K7$ & $(3,19)$-cable of $T_{2,3}$ & 6 & 21 \\
$K10$ & $(2,53)$-cable of $K1$ & 8 & 42 \\
\bottomrule
\end{tabular}
\end{table}

\subsection{Monomial normalisation}\label{sec:normalization}
The DAHA and the elliptic Hall algebra determine the invariants only up to a monomial, but the extraction of BPS invariants depends on it. We normalise by the following rules.
\begin{itemize} % tells-ok (labelled normalisation rules, referred to throughout)
\item[(N1)] The limit $t=q$ is the reduced coloured HOMFLY-PT polynomial at framing zero in the conventions of DGR.
\item[(N2)] For $\yng(1)$, the lowest $\bt$-degree in bold variables is $0$.
\item[(N2$'$)] The monomial for the colour $\lambda$ is $M_\lambda=M_\square^{|\lambda|}\,T_\lambda^{F}$, where $M_\square$ is the monomial for $\yng(1)$ fixed by (N1)(N2)
and $F=mn$ both for a torus knot $T_{m,n}$ and for an outermost cabling $C_{m,n}$.
\end{itemize}
For torus knots these rules reproduce the $27$ values of KN in all coefficients.
Rule (N2) fixes only a power of $(q/t)$ per box, which by Lemma~\ref{lem:adams} does not affect positivity.
For a coloured invariant the monomial is $\pm a^\alpha q^\beta t^\gamma$; the limit $t=q$ determines the sign, $\alpha$ and $\beta+\gamma$ (for the coloured cases (N1) then serves as a check, which we carried out).
What $t=q$ does not determine is a power $(q/t)^{c_\lambda}$ for each colour.

% [I4]
\begin{hypothesis}\label{hyp:N2}
For cables, the colour-dependent monomial $(q/t)^{c_\lambda}$ not determined at $t=q$ has the form {\rm(N2$'$)}:
it consists of a part proportional to $|\lambda|$, irrelevant for positivity, and of the power $T_\lambda^F$ of the refined twist eigenvalue.
\end{hypothesis}
The evidence is threefold:
(i) the hypothesis follows from the refined satellite formula (Hypothesis~\ref{hyp:sat});
(ii) for the non-algebraic cable $K5c$, within the one-parameter family obtained by multiplying the $\yng(2)$ invariant by $(q/t)^u$, only the value of $u$ given by (N2$'$) is compatible with the Laurent property and the transposition symmetry;
(iii) for nine cables and all colours with $|\lambda|\le3$, the form (N2$'$) agrees with the computed values.

% [F17]
\begin{proposition}[Normalisation as framing]\label{prop:N2-framing}
The monomial of {\rm(N2$'$)} can be written as
\[
M_\lambda=\bigl(M_\square\,(-(q/t)^{1/2})^{-F}\bigr)^{|\lambda|}\,\theta_\lambda^F .
\]
Thus the colour-dependent part of {\rm(N2$'$)} is exactly $F$ units of framing in the sense of Definition~\ref{def:framed}.
Shifting the reference framing by an even integer is equivalent to shifting $x$ by the same integer together with a uniform shift of the type of Lemma~\ref{lem:adams}.
Shifting by an odd integer without the Mari\~no--Vafa sign destroys the Laurent property (Proposition~\ref{prop:sign}).
\end{proposition}
\begin{proof}
$T_\lambda=\theta_\lambda\,(-(q/t)^{1/2})^{-|\lambda|}$ by Definition~\ref{def:framed}.
\end{proof}
Hypothesis~\ref{hyp:N2} can therefore be rephrased: the colour-dependent part of the normalisation is a framing by the Seifert framing $F$.
In particular, a positivity statement for all $x\in\Z$ does not depend on where $x=0$ is placed.
The remark of Aganagic and Shakirov \cite{Aganagic11} that their value for the trefoil differs from the superpolynomial of DGR by ``$-6$ units of framing and $\sqrt{q/t}$'' matches $F=mn=6$ in absolute value;
the sign depends on the choice of the positive direction of framing.

% [I5]
\begin{compresult}\label{cr:M}
For all $22$ knots computed ($9$ torus knots and $13$ cables), the monomial for $\yng(1)$ fixed by {\rm(N1)(N2)} is
\[
M_\square=(-q/t)^{b-1}\bigl(a(q/t)^{1/2}\bigr)^{\delta},
\]
where $b$ is the number of strands and $\delta$ is determined recursively by $\delta\mapsto m\delta+(m-1)(n-1)/2$ for each cabling $C_{m,n}$, starting from $\delta=(m-1)(n-1)/2$ for $T_{m,n}$.
\end{compresult}

% [I6]
\begin{hypothesis}[Refined satellite formula]\label{hyp:sat}
$\rCS_\lambda(C_{m,n}K)=\sum_\mu C_{\lambda\mu}\,\rCS_\mu(K)$, where the coefficients are the expansion coefficients of the torus knot $T_{m,n}$,
regarded as the $(m,n)$-cable of the unknot, normalised by {\rm(N2$'$)}: $C_{\lambda\mu}=M_{T_{m,n}}^{|\lambda|}\,T_\lambda^{mn}\,c_{\lambda\mu}(n)$.
\end{hypothesis}
At $t=q$ this reduces to the classical satellite formula and is compatible with (N1).
The closed formula of Computational Result~\ref{cr:M} can be derived from Hypothesis~\ref{hyp:sat} and two identities that we verified by computation.
\begin{compresult}\label{cr:tau}
Let $\tau:P_\mu\mapsto T_\mu P_\mu$ be the framing operator. For an element of slope $(m,n)$ of the elliptic Hall algebra and an integer $j$,
$\tau^jP_\lambda[p_k\mapsto P_{k(m,n)}]\cdot1=P_\lambda[p_k\mapsto P_{k(m,n+jm)}]\cdot1$.
This was checked exactly at rational points for five slopes and $j=1,-1,2$.
\end{compresult}
\begin{compresult}\label{cr:mellit}
The Poincar\'e polynomial $\mathcal{P}_{m,n}$ of the HOMFLY-PT homology of the torus knot in the Dyck-path form of Mellit \cite[Thm.~3.1]{Mellit17} satisfies
\[
\mathcal{P}_{m,n}=-\bigl(a(qt)^{-1/2}\bigr)^{\delta}\frac{1}{1-t}\,\frac{(Q_{m,n}\cdot1)[a/q-1]}{a/q-1},
\]
where $Q_{m,n}$ is the operator of Bergeron, Garsia, Leven and Xin \cite{Bergeron14}.
This was checked exactly for $(m,n)=(2,3),(3,2),(2,5),(3,4),(4,3),(3,5),(5,3)$.
\end{compresult}
\begin{proposition}\label{thm:M}
Assume Hypothesis~\ref{hyp:sat} and the identities of Computational Results~\ref{cr:tau} and~\ref{cr:mellit} for all slopes.
Then the closed formula of Computational Result~\ref{cr:M} and the form {\rm(N2$'$)} hold.
\end{proposition}
\begin{proof}[Sketch of proof]
For the cabling step, Computational Result~\ref{cr:tau} says that the framing operator maps the slope $(m,n)$ to $(m,n+jm)$.
Together with Hypothesis~\ref{hyp:sat} this gives $M_{C_{m,n}K}=M_K^{\,m}\,M_{T_{m,n}}$ and $F=mn$.
For torus knots four facts are used:
(I) the lowest coefficient in $a$ of the unnormalised value of Gorsky and Negu\c{t} is a rational $q,t$-Catalan number, which was conjectured in \cite{Gorsky13b} and follows from the theorem of Mellit \cite{Mellit16};
(II) the lowest $a$-degree is given by the theorem of Franks and Williams for positive braids;
(III) the sign and the power of $q$ at $t=q$ follow from the symmetry $q\leftrightarrow q^{-1}$;
(IV) the exponent of $(q/t)$ follows from the Dyck-path form of Mellit (Computational Result~\ref{cr:mellit}).
Details are given in Appendix~\ref{app:normalization}.
\end{proof}
\begin{remark}[Relation to the tilde normalisation]\label{rem:tilde}
Cherednik and Danilenko \cite{Cherednik14} normalise the DAHA superpolynomial by dividing by the lowest power of $a$ and by a monomial in $q,t$ so that the constant term of the part of $a$-degree zero is $1$ (tilde normalisation);
their standard variables $(q_{st},t_{st},a_{st})$ are the bold variables.
For the $19$ computed values with non-negative coefficients ($9$ torus knots and $10$ cables), our $P=\rCS_{\yng(1)}$ in bold variables equals
$\ba^{2\delta}\bq^{-2\delta}$ times the tilde-normalised polynomial.
All these $P$ are $\sigma$-invariant with even exponents of $\bq$.
\end{remark}

\subsection{Knots with homological input}
For the $801$ prime knots with at most $11$ crossings we take as the input for $\yng(1)$ the Poincar\'e polynomial of reduced HOMFLY-PT homology computed by Nakagane and Sano \cite{Nakagane21}.
The conversion to the conventions of DGR was checked by $\sigma$-invariance, by the HOMFLY-PT polynomial at $\bt=-1$, and by agreement with $\rCS_{\yng(1)}$ for torus knots.

\subsection{Conjectural formulae for two boxes}\label{sec:conj-inputs}
For hyperbolic knots, refined invariants for $\yng(2)$ and $\yng(1,1)$ are not available from a definition, and we use conjectural formulae from the literature.

For the double twist knots $K_{m,n}$ and the twist knots $K_p=K_{p,1}$
we use the formulae of KNTZ \cite{Kameyama19} for the Poincar\'e polynomials of HOMFLY-PT homology coloured by rectangular Young diagrams (refined version, conjectural).
The knot $K_{m,n}$ of KNTZ is the $(2m-1,1,2n-1)$-pretzel knot of GNSSS \cite{Gukov15}.
At $t=q$ the KNTZ formulae are proved by Fang and Zhou \cite{Fang26}.
The conversion to the $t_c$-grading of KN uses the relation of gradings of Gorsky, Gukov and Sto\v{s}i\'c \cite{Gorsky13}, $\mathbf{Q}\mapsto\bq^R$, $\mathbf{t}_r\mapsto\bq^{-1}$, $\mathbf{t}_c\mapsto\bt\bq$,
where $R$ is the number of rows of the rectangle. It was calibrated by the agreement, in all coefficients, of the trefoil $K_{1,1}=T_{2,3}$ coloured by $\yng(1)$, $\yng(2)$, $\yng(1,1)$, $\yng(2,2)$ with the elliptic Hall algebra values (see also \cite{Nawata12}).

For the thick knots $\mathbf{10_{139}}$ and $\mathbf{10_{152}}$
we use the quiver form of Sto\v{s}i\'c \cite{Stosic24} and the conjectural formula of Kucharski, Reineke, Sto\v{s}i\'c and Su\l{}kowski (KRSS) \cite{Kucharski17b} to construct the $t_c$-grading.
For $\mathbf{8_{19}}=T_{3,4}$ this procedure reproduces the elliptic Hall algebra values for all colours from $\yng(2)$ to $\yng(1,1,1)$.
\begin{compresult}\label{cr:kntz}
The refined formulae of KNTZ and the quadruply-graded formulae of GNSSS agree in all coefficients for $K_{m,n}$ with $\yng(1)$, $\yng(2)$, and, through the mirror symmetry $\mathbf{t}_r\leftrightarrow\mathbf{t}_c$, $\yng(1,1)$,
for $|m|,|n|\in\{1,2,3\}$ in each of the four quadrants of signs ($144$ cases).
Our implementation of the KNTZ formulae agrees with the values in the ancillary files of KNTZ.
\end{compresult}
\begin{corollary}
The KNTZ and GNSSS formulae for $\yng(1)$, $\yng(2)$, $\yng(1,1)$ agree for all $(m,n)$.
\end{corollary}
\begin{proof}
In each quadrant both formulae are exponential polynomials in $m$ and in $n$ with the same set of at most three bases in each variable, so agreement on a $3\times3$ grid implies agreement everywhere.
\end{proof}
Statements with these inputs are statements of positivity \emph{given the conjectural formulae}; they are of a different nature from the three-box results of \S\ref{sec:three}, whose inputs are defined.

\subsection{Exact computation}\label{sec:exact}
$\brCS_\lambda$ and $\hf_\rho$ are obtained by evaluation modulo $31$-bit primes followed by interpolation in each variable, not by symbolic manipulation.
When coefficients are large, two or three primes and the Chinese remainder theorem are used.
Every polynomial obtained is checked at points not used in the interpolation.
The largest example, $K10$ with three boxes, has about $1.83$ million terms in the spin form, with coefficients up to about $1.1\times10^{19}$.
Running times and memory are given in Appendix~\ref{app:computation}.
